%----------------------------------------------------------------------------------------------------------------------------
\chapter*{Epilogo}
\addcontentsline{toc}{chapter}{Epilogo}
%----------------------------------------------------------------------------------------------------------------------------

È il momento di fare il punto. Non nel senso di un semplice inventario — i sommari di ciascun capitolo hanno già assunto questo compito — ma nel senso di guardare indietro all'intero percorso per vederne la forma, capire fin dove arriva e riconoscere con onestà dove si ferma.


%----------------------------------------------------------------------------------------------------------------------------
\section*{Ciò che è stato costruito}
%----------------------------------------------------------------------------------------------------------------------------

Il filo conduttore di questo volume è stato il \textit{corpo rigido}: un modello che rinuncia deliberatamente a descrivere la deformazione dei materiali per concentrarsi sulla struttura logica della meccanica. La scelta non è una semplificazione ingenua: è una strategia consapevole. Privando la teoria del legame costitutivo --- il terzo dei tre pilastri introdotti nell'Introduzione --- si mette in evidenza la forza dei primi due.

E la forza emerge con chiarezza. Le equazioni di congruenza e le equazioni di equilibrio, prese da sole, bastano a descrivere completamente il comportamento di qualsiasi sistema isostatico: una trave a mensola, un arco a tre cerniere, una travatura reticolare isostatica, un ``telaio'' con il numero giusto di vincoli. In tutti questi casi il problema cinematico ammette soluzione unica, il problema statico ammette soluzione unica, e le due soluzioni sono collegate dal principio dei lavori virtuali senza che mai compaia una proprietà del materiale. Il legame costitutivo è assente, eppure la teoria funziona.

Questo risultato non è banale. Significa che la dualità cinematica-statica --- il filo conduttore di tutti e tre i volumi --- non è una proprietà del materiale ma una proprietà della struttura logica della meccanica. L'operatore cinematico $\bm{A}$ e il suo trasposto $\bm{A}^\top$, che è l'operatore di equilibrio, sono legati non da un'ipotesi fisica ma dall'identità bilineare del lavoro virtuale. Questa connessione è algebrica prima ancora che meccanica, e il suo significato è stato reso visibile dall'analisi dei quattro sottospazi fondamentali: meccanismi e stati di autoequilibrio, cedimenti compatibili e forze equilibrabili, si collocano in una struttura geometrica che non dipende dal materiale ma soltanto dalla topologia del sistema e dalla geometria dei vincoli.


La seconda conquista del volume è la \textit{classificazione}. Il conteggio $\nu = 3n_c - m$ --- gradi di libertà nominali meno molteplicità vincolare --- fornisce la prima informazione qualitativa, ma il rango della matrice $\bm{A}$ è ciò che determina davvero la natura del sistema. Quattro categorie: \textit{isostatico} ($\nu = 0$, $\bm{A}$ a rango pieno), \textit{labile} ($\nu > 0$, oppure $\nu \leq 0$ con difetto di rango), \textit{iperstatico} ($\nu < 0$, oppure $\nu \geq 0$ con difetto di rango), \textit{degenere} (labile e iperstatico simultaneamente). Questa classificazione vale per il sistema di corpi rigidi generico e si trasferisce senza modifiche alle travature reticolari e ai telai piani: la struttura algebrica è la stessa, cambiano soltanto le dimensioni delle matrici.

La terza conquista è la \textit{specializzazione alla trave}. Introdotto il riferimento intrinseco e le caratteristiche della sollecitazione --- forza normale $N$, taglio $T$, momento flettente $M$ --- la teoria del corpo rigido diventa uno strumento operativo preciso: i diagrammi delle caratteristiche della sollecitazione, costruiti con il metodo per equivalenza e controllati con le equazioni locali di equilibrio (indefinite e di salto), permettono di descrivere compiutamente lo stato interno di qualsiasi sistema isostatico di travi.

\begin{oss}
	La trattazione sviluppata in questo volume è riferita, salvo brevi cenni, ai \textit{sistemi piani}: corpi rigidi che si muovono in un piano, travi con asse rettilineo nel piano, vincoli e carichi contenuti nello stesso piano. Questa scelta è dettata da ragioni didattiche precise: il caso piano riduce la complessità notazionale, rende visibili le costruzioni geometriche e consente di concentrarsi sulla struttura logica senza che l'algebra vettoriale tridimensionale oscuri il ragionamento meccanico. L'estensione al caso spaziale --- sistemi di corpi rigidi nello spazio, travi curve, telai tridimensionali --- non introduce complessità concettuali aggiuntive: la struttura delle equazioni, la dualità fra operatore cinematico e operatore di equilibrio, la classificazione in base al rango della matrice $\bm{A}$ sono identiche. Cambiano le dimensioni delle matrici e il numero dei gradi di libertà per corpo (sei nello spazio invece di tre nel piano), ma non la grammatica.
\end{oss}

%----------------------------------------------------------------------------------------------------------------------------
\section*{I limiti del modello rigido}
%----------------------------------------------------------------------------------------------------------------------------

Il modello del corpo rigido, nella sua eleganza, porta in sé due limitazioni simmetriche. Riconoscerle non è ammettere una debolezza: è il passo che apre le porte ai volumi successivi.

\paragraph{I sistemi labili.}
Un sistema labile possiede gradi di labilità $g_\ell > 0$: ammette spostamenti rigidi non impediti dai vincoli. Il problema cinematico ha infinite soluzioni --- il sistema è un meccanismo, capace di muoversi senza deformarsi. La controparte statica è immediata: il problema di equilibrio $\bm{A}^\top \bm{r} = -\bm{f}$ è sovradeterminato, e per un carico $\bm{f}$ generico non ammette soluzione. Un sistema labile non può essere in equilibrio sotto forze arbitrarie: esistono carichi che non può sostenere.

Questa limitazione non è necessariamente un difetto da eliminare, ma piuttosto un segnale che il modello deve essere arricchito. Mantenendo l'ipotesi di rigidità degli elementi e introducendo organi a \textit{deformabilità concentrata} --- molle, cuscinetti elastici, giunti a rigidezza finita --- nei punti in cui il sistema si muove, i meccanismi acquistano rigidezza e l'equilibrio sotto carichi generici torna possibile. Se gli organi sono in numero pari ai gradi di labilità, e ben posti --- ciascuno a contrastare un meccanismo distinto --- il sistema diventa staticamente determinato: le forze negli organi seguono dal solo equilibrio, e il legame costitutivo interviene unicamente a recuperare gli spostamenti. Esso torna necessario anche per le forze soltanto quando gli organi eccedono i gradi di labilità, o rispondono in modo non lineare --- elastico-plastico. Su questo terreno, il problema elastico ed elasto-plastico di sistemi rigidi a deformabilità concentrata, si muove il testo \cite{LPDN}, a cui si rimanda per una trattazione completa.


\paragraph{I sistemi iperstatici.}
Un sistema iperstatico possiede gradi di iperstaticità $g_\imath > 0$: le equazioni di equilibrio sono in numero inferiore alle  reazioni vincolari, e il problema statico ammette infinite soluzioni. Esistono stati di autoequilibrio --- distribuzioni di reazioni vincolari compatibili con forza esterna nulla --- che il solo equilibrio non può determinare.

La controparte cinematica è altrettanto netta: il problema cinematico $\bm{A}\bm{u} = \bm{s}$ è sovradeterminato. Per cedimenti $\bm{s}$ generici, non esistono spostamenti rigidi compatibili. Il sistema iperstatico non può accomodare cedimenti vincolari arbitrari senza deformarsi. Ed è esattamente qui che il modello rigido si arresta: per determinare quale delle infinite soluzioni statiche sia quella effettivamente realizzata, e per rendere compatibili i cedimenti vincolari, occorre sapere \textit{come si deforma} la struttura. Occorre, in altre parole, il legame costitutivo.

I due limiti sono perfettamente duali: la labilità è indeterminazione cinematica che produce impossibilità statica; l'iperstaticità è indeterminazione statica che produce impossibilità cinematica. A ciascuna corrisponde una risposta, e le due risposte sono anch'esse duali nel loro spirito, se non nelle loro conseguenze tecniche.


%----------------------------------------------------------------------------------------------------------------------------
\section*{Il terzo pilastro}
%----------------------------------------------------------------------------------------------------------------------------

I due limiti condividono una radice e un completamento. La radice è la stessa --- un nucleo non banale: i meccanismi in $\Nc(\bm{A})$, gli stati di autoequilibrio in $\Nc(\bm{A}^\top)$ --- e a completarli è, in entrambi i casi, il terzo pilastro che il modello rigido aveva deliberatamente messo da parte: il legame costitutivo. Ma il pilastro entra nei due problemi a profondità diverse, e la ragione sta nel modo opposto in cui l'equilibrio fallisce. Nel sistema labile l'equilibrio è \textit{sovradeterminato}: le reazioni incognite sono troppo poche per le equazioni, e per un carico generico non esiste soluzione. Nel sistema iperstatico è \textit{sottodeterminato}: le incognite sono troppe, e le soluzioni sono infinite. A un difetto di incognite e a un loro eccesso corrispondono due rimedi opposti.

Un difetto di incognite si corregge aggiungendone. Per curare la labilità non occorre toccare la rigidità degli elementi: un meccanismo è un moto che i corpi compiono restando rigidi, e lo si contrasta agendo sulle connessioni, dove il moto relativo ha luogo. Una connessione è un oggetto discreto, e la forza che vi nasce è una nuova incognita altrettanto discreta. Bastano $g_\ell$ organi a \textit{deformabilità concentrata}, ben posti, perché le incognite tornino pari alle equazioni: il sistema diventa staticamente determinato e le forze negli organi seguono dal solo equilibrio. Il legame costitutivo non serve a trovarle --- serve a recuperare gli spostamenti, e a governare le forze stesse solo quando gli organi eccedono i gradi di labilità o rispondono in modo non lineare. È la strada del testo \cite{LPDN}: il problema elastico ed elasto-plastico formulato direttamente sui gradi di libertà già presenti nel modello discreto.

Un eccesso di incognite chiede il rimedio opposto. Non si corregge aggiungendone altre, ma aggiungendo \textit{equazioni}; e le sole equazioni nuove disponibili sono quelle di congruenza. Scriverle, però, richiede di sapere come la struttura si deforma --- cioè il legame costitutivo --- che entra così nella determinazione delle forze stesse, e non a valle di esse. E qui la deformabilità non può essere concentrata. Il vincolo che si è costretti a rilassare è la rigidità del membro, e la rigidità non è una condizione imposta in un punto, ma in \textit{ogni} punto materiale: è un vincolo di campo. La sua grandezza duale è perciò anch'essa un campo --- deformazioni e tensioni definite in ciascun punto dell'asse --- e la congruenza che richiude il problema iperstatico non è una condizione imposta in un punto, ma il risultato della deformazione che si accumula lungo l'intero elemento, sezione dopo sezione. La deformabilità che cura l'iperstaticità è dunque \textit{diffusa}: non per convenzione, ma perché la grandezza coniugata alla rigidità è distribuita per sua natura. È la strada del Volume~II.


È questo il passaggio dal primo al secondo volume: dalla \textit{trave rigida} alla \textit{trave deformabile}. Non un cambio di paradigma, ma un completamento. La struttura logica --- congruenza, equilibrio, dualità --- rimane identica. La matrice $\bm{A}$, che esprime le condizioni di vincolo con l'esterno e fra i corpi, rimane invariata; ciò che cambia è il termine noto $\bm{s}$, che ora raccoglie non solo i cedimenti vincolari esterni ma anche il contributo della deformazione interna degli elementi. Compaiono le equazioni di congruenza interna --- quelle che assicurano la continuità del materiale, l'assenza di lacerazioni e compenetrazioni --- e il legame costitutivo che le connette alle sollecitazioni. Ma il lettore che ha percorso questo volume riconoscerà, in ogni passaggio del secondo, la stessa grammatica. Sarà una lettura di conferma, non di scoperta.




%----------------------------------------------------------------------------------------------------------------------------
\section*{Verso il Volume~II}
%----------------------------------------------------------------------------------------------------------------------------

Il Volume~II ha per protagonista la \textit{trave rettilinea deformabile}, concepita e sviluppata nella sua piena generalità: una trave immersa nello spazio euclideo tridimensionale, con sei gradi di libertà per sezione --- tre traslazioni e tre rotazioni --- e sei caratteristiche della sollecitazione corrispondenti. Il recupero della tridimensionalità non richiede di rivedere i fondamenti: la struttura logica rimane quella del Volume~I, e il lettore riconoscerà in ogni equazione la stessa dualità cinematica-statica, ora dispiegata nella sua forma più generale. Le sue equazioni si derivano dal continuo tridimensionale di Cauchy mediante le ipotesi cinematiche di riduzione dimensionale --- Eulero-Bernoulli e Timoshenko --- che trasformano il problema tridimensionale in un problema formulato sull'asse della trave. Le caratteristiche della sollecitazione $N$, $T$, $M$ --- introdotte in questo volume come reazioni dell'equilibrio --- acquistano nel secondo volume una seconda lettura: sono le variabili duali delle deformazioni generalizzate, e il loro legame con queste ultime è il legame costitutivo della trave.

Su questa base si sviluppano i metodi di analisi delle strutture iperstatiche. Il Volume~II li affronta secondo due punti di vista complementari: quello della trave come continuo monodimensionale, in cui i campi di spostamento e sollecitazione sono descritti lungo l'intero asse, e quello nodale, in cui le variabili primarie sono concentrate ai nodi di connessione. Per ciascuno dei due punti di vista si presentano il \textit{metodo delle forze} --- che sceglie le incognite iperstatiche come variabili primarie e impone la compatibilità --- e il \textit{metodo degli spostamenti} --- che sceglie gli spostamenti come variabili primarie e impone l'equilibrio. I due metodi sono duali, e la dualità che in questo volume era algebrica --- $\bm{A}$ e $\bm{A}^\top$ --- diventa nel secondo volume meccanica: forze e spostamenti, deformabilità e rigidezza.

Chi ha assimilato i contenuti di questo primo volume possiede già la mappa del territorio. Il secondo volume la percorre.